\documentclass[aps,prd,twocolumn,superscriptaddress,nofootinbib]{revtex4-2}

\usepackage{amsmath,amssymb,amsfonts}
\usepackage{graphicx}
\usepackage{hyperref}
\usepackage{bm}
\usepackage{booktabs}
\usepackage{amsthm}
\newtheorem{theorem}{Theorem}
\newtheorem{proposition}[theorem]{Proposition}
\newtheorem{conjecture}[theorem]{Conjecture}

\newcommand{\Spec}{\mathrm{Spec}}
\newcommand{\Tr}{\mathrm{Tr}}
\newcommand{\cd}{\mathrm{cd}}
\newcommand{\zetanot}{\zeta_{\neg}}
\newcommand{\KMS}{\mathrm{KMS}}

\begin{document}

\title{Spectral Arithmetic Warp:\\
Vacuum Engineering via the Bost-Connes Phase Transition\\
and Euler Product Amplification}

\author{Wright Brothers}
\affiliation{Independent Research}
\date{\today}

\begin{abstract}
We develop five approaches to manipulating quantum vacuum energy
using the arithmetic structure of the Riemann zeta function,
as realized in the Bost-Connes (BC) C$^*$-dynamical system.
(1)~\emph{Spectral action approach}: local perturbation $\delta D$ of the
Dirac operator changes geometry via Connes' spectral action principle,
with the BC heat kernel $\Tr(e^{-tH}) = \zeta(t)$ providing the bridge
between number theory and gravity.
(2)~\emph{$\beta$-pumping}: population inversion of cavity phonons
creates a transient $\beta < 0$ state (negative temperature),
realizing local dark energy ($\beta = -1$) via a standard laboratory
technique (Purcell-Pound protocol).
(3)~\emph{Euler product amplification}: removing $K$ primes from the
mode spectrum amplifies Casimir energy by $\prod_{i=1}^K (p_i^3 - 1)$;
14 primes yield $10^{48}\times$.
(4)~\emph{Inverse spectral design}: a graded material whose natural
eigenfrequencies match the prime-sieved set, physically instantiating
the Euler product as an ``arithmetic drum.''
(5)~\emph{Dirac spectral modulation}: $\pi$-phase rotation of all
Euler factors (Liouville function vacuum, $\zeta(t)/\zeta(2t)$)
yields $G_{\mathrm{eff}} = -2G$ (repulsive gravity); muting $K \geq 2$
primes yields $G_{\mathrm{eff}} = 0$ (gravity nullification).
A combined two-layer warp bubble (outer: repulsive, inner: gravity-free)
requires $\sim 10^3$\,J of positive energy---44 orders of magnitude
below Alcubierre.
We prove that rapid switching yields exactly zero net energy
(vacuum Carnot theorem), confirming energy conservation.
The decisive experimental test---comparing Casimir energy in a
2-prime-sieved cavity versus a uniform cavity (predicted ratio 182:1
vs.\ 1:3)---costs approximately \$2{,}200.
\end{abstract}

\maketitle

%=======================================================================
\section{Introduction}
%=======================================================================

The Alcubierre warp metric~\cite{Alcubierre1994} requires negative
energy density (violation of the weak energy condition), typically
estimated at $\sim 10^{47}$\,J of exotic matter.
This paper explores whether the arithmetic structure of the quantum
vacuum---specifically the Euler product decomposition of the Riemann
zeta function---offers alternative mechanisms for vacuum engineering.

Our starting point is the Bost-Connes (BC) C$^*$-dynamical
system~\cite{BostConnes1995}, whose partition function is the Riemann
zeta function:
\begin{equation}
  Z(\beta) = \Tr(e^{-\beta H}) = \sum_{n=1}^\infty n^{-\beta} = \zeta(\beta).
  \label{eq:bc_partition}
\end{equation}
The BC system exhibits a phase transition at $\beta = 1$, with
spontaneous breaking of the Galois symmetry
$\mathrm{Gal}(\mathbb{Q}^{\mathrm{ab}}/\mathbb{Q})$ for $\beta > 1$.

In previous work~\cite{WB_arithmetic_universe}, we showed that the
dark energy density $\Omega_\Lambda = 2\pi/9$ arises from the
completed zeta function at negative temperature:
\begin{equation}
  \Omega_\Lambda = \frac{d}{\cd} \,|\xi_{\neg 2}(-1)|
  = \frac{4}{3} \cdot \frac{\pi}{6} = \frac{2\pi}{9} \approx 0.698,
\end{equation}
where $d = 4$, $\cd = \cd(\Spec(\mathbb{Z})) = 3$ (Artin-Verdier),
and $\xi_{\neg 2}$ is the completed zeta with the $p = 2$ Euler
factor removed.

Here we ask: \emph{can the same arithmetic structure be exploited
for local vacuum engineering?}

%=======================================================================
\section{The $\beta$-line of physics}
\label{sec:beta_line}
%=======================================================================

The inverse temperature $\beta$ of the BC system parameterizes
distinct physical regimes:

\begin{center}
\begin{tabular}{lll}
\toprule
$\beta$ & $\zeta(\beta)$ & Physical sector \\
\midrule
$\beta > 1$ & convergent & Perturbative QFT \\
$\beta = 1$ & pole & Phase transition \\
$0 < \beta < 1$ & & Grand unification \\
$\beta = 0$ & $-1/2$ & Topological \\
$\beta = -1$ & $-1/12$ & Dark energy \\
$\beta = -3$ & $1/120$ & 3D Casimir / K-theory \\
\bottomrule
\end{tabular}
\end{center}

The functional equation $\xi(\beta) = \xi(1-\beta)$ establishes a
duality: $\beta = -1 \leftrightarrow \beta = 2$, mapping dark energy
to perturbative physics.

The sign of $\zeta_{\neg 2}$ alternates at odd negative integers:
$\zeta_{\neg 2}(-1) = +1/12$ (positive, repulsive),
$\zeta_{\neg 2}(-3) = -7/120$ (negative), establishing a
dimension-dependent sign structure.

%=======================================================================
\section{Approach 1: Spectral action}
\label{sec:spectral}
%=======================================================================

Connes' spectral action principle~\cite{Chamseddine1996} gives
\begin{equation}
  S = \Tr\bigl(f(D^2/\Lambda^2)\bigr)
  \sim \sum_k f_k \Lambda^k a_k(D^2),
\end{equation}
where $a_k$ are the Seeley-DeWitt coefficients encoding geometry.

For $D^2 = H_{\mathrm{BC}}$ with heat kernel~(\ref{eq:bc_partition}):
$S(\Lambda) = \zeta(1/\Lambda^2)$.
The pole at $\Lambda = 1$ ($\leftrightarrow \beta = 1$) is the
BC phase transition, where the spectral action diverges.

A local perturbation $D \to D + V(x)$ modifies the spectral action
density:
\begin{equation}
  s(x) = \zeta(1/\Lambda^2) \, e^{-V(x)/\Lambda^2},
\end{equation}
producing expansion ($V < 0$) or contraction ($V > 0$)---an
Alcubierre-like profile in spectral language.

\begin{proposition}[Finite energy, divergent response]
Near the BC phase transition $\beta \to 1^+$: the energy cost
$E(\beta) = -\zeta'(\beta)/\zeta(\beta)$ remains finite,
while the geometric response $\zeta(\beta) \sim 1/(\beta-1)$
diverges.
\end{proposition}

However, the specific heat $C(\beta)$ also diverges as $1/(\beta-1)^2$,
limiting the leverage ratio $\chi/C$ to $O(1)$.
The product $\delta g \times E_{\mathrm{wall}}$ scales identically
to general relativity, confirming that the spectral approach does not
circumvent energy-geometry proportionality by itself.

%=======================================================================
\section{Approach 2: $\beta$-pumping}
\label{sec:pumping}
%=======================================================================

The $\beta = -1$ state (dark energy) can be created locally via
\emph{population inversion}---a standard technique in NMR and
quantum optics (Purcell-Pound, 1951).

\textbf{Protocol.}
(1)~Prepare a BAW crystal in DD thermal equilibrium at temperature $T$.
(2)~Rapidly switch one boundary from D to N (switching time
$\tau < 1/(2\omega_0)$).
(3)~The phonon populations are now inverted relative to the new
mode spectrum: $\beta_{\mathrm{eff}} < 0$.
(4)~During the relaxation time $T_1$, the cavity exhibits negative
effective pressure.

The energy density of the population-inverted state is
$\rho_{\mathrm{dyn}} \sim N_{\mathrm{modes}} \, k_B T / V$,
which exceeds the static Casimir energy $\rho_{\mathrm{Cas}}
\sim \hbar c / d^4$ by a factor $\sim \bar{n} \times N_{\mathrm{modes}}
\sim 10^{20}$.

The effect is \emph{transient} (duration $\sim T_1$) but
\emph{repeatable} (pulsed operation at $\sim 1/(2T_1)$).
This constitutes a propellantless acoustic thruster concept,
though with small forces ($\sim 10^{-13}$\,N per crystal).

%=======================================================================
\section{Approach 3: Euler product amplification}
\label{sec:euler}
%=======================================================================

The Euler product
$\zeta(s) = \prod_p (1 - p^{-s})^{-1}$
decomposes the vacuum into independent prime channels.
Removing primes $p_1, \ldots, p_K$ gives
\begin{equation}
  \zeta_{\neg\{p_i\}}(s) = \prod_{i=1}^K (1 - p_i^{-s}) \, \zeta(s).
  \label{eq:euler_muted}
\end{equation}

At $s = -3$ (3D Casimir), each removed prime amplifies the spectral
zeta by $|1 - p^3| = p^3 - 1$:

\begin{center}
\begin{tabular}{rrl}
\toprule
$K$ & Primes & Amplification \\
\midrule
1 & $\{2\}$ & $7\times$ \\
2 & $\{2,3\}$ & $182\times$ \\
4 & $\{2,3,5,7\}$ & $7.7 \times 10^6$ \\
14 & $\{2,\ldots,43\}$ & $1.9 \times 10^{48}$ \\
\bottomrule
\end{tabular}
\end{center}

For even $K$, the Casimir energy is negative (exotic matter).
At $K = 14$, the amplification exceeds the Alcubierre requirement.

\subsection{Safety analysis}

Previous work identified that removing primes from the \emph{global}
vacuum changes L-function values linked to coupling constants
($\theta_W$, $\alpha_s$).
However, \emph{local} mode filtering in a cavity---analogous to
the Casimir effect between plates, which modifies the mode spectrum
without changing $\alpha$---may not affect global coupling constants.

This distinction is \emph{not rigorously established}.
The BC framework suggests L-function values change under
Euler factor removal, while standard QFT UV/IR decoupling
suggests they do not.

\begin{conjecture}[Regulator conjecture]
For a physical cavity whose spectral zeta function equals
$\zeta_{\neg\{p_i\}}(s)$, the Casimir energy equals the
zeta-regularized value $-({\hbar c}/{d^3}) \times
\zeta_{\neg\{p_i\}}(-3) \times (\text{geometric factor})$.
\end{conjecture}

This conjecture is experimentally testable (Section~\ref{sec:experiment}).

\subsection{Vacuum Carnot theorem}

\begin{theorem}[Energy conservation under rapid switching]
A cycle of boundary-condition switching
$\mathrm{DD} \to \mathrm{DN}_{\neg 2,3} \to \mathrm{DD}$
with intermediate mechanical work yields exactly zero net energy:
\begin{equation}
  \oint \delta E = 0
\end{equation}
for any path in $(L, \text{BC})$ parameter space.
\end{theorem}

\begin{proof}
The total energy is $E(L, \text{BC}) = -A_{\text{BC}} / L^3$
where $A_{\text{BC}}$ depends only on the boundary condition.
In a cycle: $\Delta E_{\text{switch}} + W_{\text{mech}} = 0$ at
each step, and the sum telescopes to zero.
\end{proof}

This rules out energy extraction from vacuum switching
and confirms the theoretical self-consistency of the framework.

%=======================================================================
\section{Approach 4: Inverse spectral design}
\label{sec:inverse}
%=======================================================================

Rather than suppressing modes with bandgap filters (\emph{subtractive}),
we propose designing a graded material whose \emph{natural}
eigenfrequencies match the prime-sieved set (\emph{constructive}).

The inverse Sturm-Liouville problem: given target eigenvalues
$\lambda_k = m_k^2$ where $\{m_k\}$ is the prime-sieved sequence,
find a potential $V(x)$ on $[0, \pi]$ such that
\begin{equation}
  -y'' + V(x) y = \lambda y, \quad y(0) = y(\pi) = 0
\end{equation}
has spectrum $\{\lambda_k\}$.

By the Gel'fand-Levitan theorem, $V(x)$ exists and is unique
(given appropriate normalizing constants).
The potential encodes the Eratosthenes sieve in its shape.

\textbf{Material realization.}
$V(x)$ maps to position-dependent stiffness $c(x)^2$ in an acoustic
resonator. This can be fabricated as a compositionally graded
piezoelectric thin film (e.g., Al$_x$Ga$_{1-x}$N with programmed $x(z)$)
via molecular beam epitaxy or atomic layer deposition.

The \emph{arithmetic drum}---a single graded crystal---replaces the
14-period phononic sieve, with no bandgap leakage and exact eigenvalues
by construction.

%=======================================================================
\section{Approach 5: Dirac spectral modulation}
\label{sec:dirac}
%=======================================================================

Rather than removing Euler factors (muting) or designing specific
eigenfrequencies (inverse spectral), we consider \emph{phase-rotating}
the Euler factors: replacing $(1 - p^{-t})^{-1}$ with
$(1 - e^{i\phi} p^{-t})^{-1}$.

\subsection{Gravity nullification ($G_{\mathrm{eff}} = 0$)}

The Seeley-DeWitt coefficient $a_2$ (which determines $G$) is
proportional to the heat kernel derivative $K'(0)$.
When $K$ primes are muted, $\zeta_{\neg\{p_i\}}(t)$ has a $K$-fold
zero at $t = 0$:
\begin{equation}
  \zeta_{\neg\{p_i\}}(t) = \prod_{i=1}^K (1 - p_i^{-t}) \, \zeta(t)
  \sim t^K \quad (t \to 0).
\end{equation}

For $d = 4$, the heat kernel expansion
$K(t) = a_0 t^{-2} + a_2 t^{-1} + a_4 + \cdots$
requires $a_0 = a_2 = 0$ when $K \geq 2$.

\begin{proposition}[Gravity nullification]
If the spectral zeta of a local region equals
$\zeta_{\neg\{p_1, p_2\}}(t)$ for any two primes $p_1, p_2$,
then $a_0 = a_2 = 0$ in that region, implying
$G_{\mathrm{eff}} = 0$ and no cosmological constant.
\end{proposition}

\subsection{Anti-gravity ($G_{\mathrm{eff}} < 0$) via Liouville rotation}

Setting $\phi = \pi$ for \emph{all} primes replaces each
$(1 - p^{-t})^{-1}$ with $(1 + p^{-t})^{-1}$.
Using the identity
\begin{equation}
  \prod_p (1 + p^{-t})^{-1} = \frac{\zeta(2t)}{\zeta(t)},
  \label{eq:liouville}
\end{equation}
which follows from $\zeta(t)/\zeta(2t) = \prod_p (1 + p^{-t})$,
the all-rotated heat kernel is $K_\pi(t) = \zeta(t)/\zeta(2t)$.

At $t = 0$: $K_\pi(0) = \zeta(0)/\zeta(0) = 1$ (no zero).
The derivative:
\begin{equation}
  K_\pi'(0) = \frac{d}{dt}\biggl[\frac{\zeta(t)}{\zeta(2t)}\biggr]_{t=0}
  = -\frac{\zeta'(0)}{\zeta(0)} = \log(2\pi).
\end{equation}

Compare the normal case $K'(0) = \zeta'(0) = -\tfrac{1}{2}\log(2\pi)$.
The ratio
\begin{equation}
  \frac{G_{\mathrm{eff}}}{G} = \frac{K_\pi'(0)}{K'(0)}
  = \frac{\log(2\pi)}{-\tfrac{1}{2}\log(2\pi)} = -2.
\end{equation}

\begin{proposition}[Anti-gravity]
If all prime Euler factors are $\pi$-rotated
(Liouville function vacuum), the effective gravitational coupling
reverses: $G_{\mathrm{eff}} = -2G$.
This is repulsive gravity.
\end{proposition}

\textbf{Implementation.}
Each transmon qubit (tuned to $\omega_p = \log p / \tau$) applies
a $Z$-gate ($\pi$-phase shift) rather than absorbing the mode.
The $Z$-gate is a standard single-qubit operation in circuit QED,
requiring $\sim 10^{-20}$\,J per gate.

\subsection{The combined arithmetic warp bubble}

A two-layer structure achieves both repulsion and decoupling:
\begin{itemize}
\item \textbf{Outer shell:} all primes $\pi$-rotated $\to$
  $G_{\mathrm{eff}} = -2G$ (repulsive, pushes surrounding space).
\item \textbf{Inner region:} 2 primes muted $\to$
  $G_{\mathrm{eff}} = 0$ (ship floats, no tidal forces).
\item \textbf{Transition wall:} metamaterial boundary between layers.
\end{itemize}

Total energy: $\sim 10^3$\,J (wall structure),
compared to $\sim 10^{47}$\,J of exotic matter for
standard Alcubierre.

\subsection{Computational constraints inside the bubble}

Muting $p = 2$ weakens binary arithmetic (parity-based error correction
loses its algebraic foundation in $\mathbb{F}_2$).
Inside a $K$-prime-muted bubble, computation must use a number base
$b > p_K$ (the largest muted prime):

\begin{center}
\begin{tabular}{ccc}
\toprule
Primes muted & Largest $p$ & Safe base \\
\midrule
$\{2,3\}$ & 3 & 5 (quinary) \\
$\{2,\ldots,13\}$ & 13 & 17 \\
$\{2,\ldots,43\}$ & 43 & 47 \\
\bottomrule
\end{tabular}
\end{center}

Error correction over $\mathrm{GF}(47)$ or larger is standard
in cryptographic applications (Reed-Solomon codes).

\subsection{Entropy maintenance cost}

The Landauer bound gives the minimum energy to maintain the
muted state:
$W \geq k_B T \,|\Delta S|$ per mode,
where $\Delta S \approx 0.39$ (from BC entropy calculation).
At 10\,mK: $W \sim 10^{-8}$\,J\,m$^{-3}$---negligible compared
to the wall energy.

%=======================================================================
\section{The decisive experiment}
\label{sec:experiment}
%=======================================================================

The regulator conjecture (Section~\ref{sec:euler}) makes a sharp,
falsifiable prediction.

\textbf{Setup.} Three BAW crystals:
\begin{itemize}
\item[A.] Uniform (DD): eigenvalues $\propto n^2$ for all $n$.
\item[B.] DN boundary: eigenvalues $\propto (2n-1)^2$ (odd $n$ only).
  Implements $p = 2$ muting.
\item[C.] DN + internal Bragg reflectors at $L/3$, $2L/3$:
  eigenvalues $\propto m^2$ for $m$ not divisible by 2 or 3.
  Implements $p = 2, 3$ muting.
\end{itemize}

\textbf{Measurement.} Temperature-sweep differential Casimir energy
extraction: $\Delta E = E(T_2) - E(T_1)$ for each crystal.

\textbf{Prediction.}
\begin{center}
\begin{tabular}{lcc}
\toprule
& Euler product & Mode counting \\
\midrule
$|E_C / E_A|$ & 182 & $1/3$ \\
\bottomrule
\end{tabular}
\end{center}

A factor of 546 separates the two predictions.
This is easily distinguishable experimentally.

\textbf{Internal Bragg reflectors.}
Mo/AlN Bragg stacks (5 pairs, 500\,nm total) deposited at positions
$L/3$ and $2L/3$ within a 142\,$\mu$m quartz crystal.
The acoustic impedance contrast (quartz: 15.2\,MRayls,
Mo: 63.1\,MRayls) gives $> 99\%$ reflectivity per stack.

\textbf{Cost:} \$2{,}200 total (including custom crystal fabrication).
\textbf{Timeline:} 4--6 months.

%=======================================================================
\section{Discussion}
%=======================================================================

\subsection{What the arithmetic structure provides}

The four approaches share a common insight: the Euler product
decomposition of $\zeta(s)$ makes the vacuum energy \emph{multiplicatively
decomposable} into independent prime channels.
This is a structure that does not exist in standard QFT
(where the vacuum is treated as a single entity).

If the regulator conjecture is confirmed experimentally,
the implications extend beyond vacuum engineering:
\begin{itemize}
\item Dark energy ($\Omega_\Lambda = 2\pi/9$) and Casimir energy
  are different evaluations of the \emph{same} L-function $\zeta_{\neg 2}$.
\item The Euler product is not merely a mathematical identity
  but a physical decomposition of the vacuum into prime sectors.
\item Each prime sector can be independently manipulated,
  with multiplicative effects on vacuum energy.
\end{itemize}

\subsection{What remains unresolved}

\begin{enumerate}
\item \textbf{Local vs.\ global muting:}
  Does local mode filtering change coupling constants?
  The BC framework suggests yes (L-function values change);
  standard QFT UV/IR separation suggests no.
  The experiment decides.

\item \textbf{Energy scaling:}
  The product $\delta g \times E_{\mathrm{wall}} \sim \text{const}$
  holds in both GR and the spectral approach.
  The BC phase transition offers some leverage but not unlimited.
  Practical warp remains energy-limited.

\item \textbf{Regulator conjecture:}
  Whether $\zeta_{\mathrm{muted}}(-3)$ faithfully gives the
  physical Casimir energy for highly non-standard mode spectra
  is the central open question.
\end{enumerate}

\subsection{Regardless of warp}

Even if the framework does not lead to practical warp drive,
the experimental program has independent value:
\begin{itemize}
\item Verification of the arithmetic vacuum structure
  ($\zeta_{\neg 2}$ as physical vacuum).
\item A new test of zeta-function regularization
  (the most widely used but least tested regularization in QFT).
\item Connection between laboratory Casimir physics and
  cosmological dark energy via the same L-function.
\item A novel propellantless thruster concept ($\beta$-pumping).
\end{itemize}

%=======================================================================
\section{Conclusion}
%=======================================================================

We have shown that the Bost-Connes C$^*$-dynamical system provides
a unified framework connecting dark energy, Casimir physics, and
vacuum engineering through the arithmetic of the Riemann zeta function.

Five distinct approaches to vacuum manipulation were developed:
the spectral action approach provides conceptual clarity,
$\beta$-pumping offers immediate experimental access,
Euler product amplification gives the largest potential effect,
inverse spectral design provides the cleanest implementation,
and Dirac spectral modulation yields the most radical prediction---local
control of the effective gravitational coupling, including
$G_{\mathrm{eff}} = 0$ (gravity nullification via 2-prime muting) and
$G_{\mathrm{eff}} = -2G$ (anti-gravity via all-prime $\pi$-rotation
through the Liouville function).

The combined arithmetic warp bubble---an outer repulsive shell
($G = -2G$) surrounding a gravity-free interior ($G = 0$)---requires
$\sim 10^3$\,J of positive energy, 44 orders of magnitude less than
Alcubierre's $10^{47}$\,J of exotic matter.

All five approaches converge on a single decisive experiment: measuring the
Casimir energy ratio in a 2-prime-sieved acoustic cavity.
The predicted ratio (182:1 vs.\ 1:3) is large enough for unambiguous
determination with modest equipment (\$2{,}200, 4--6 months).

If confirmed, the result would establish the Euler product as a
physical---not merely mathematical---decomposition of the quantum
vacuum, opening a path from number theory to vacuum engineering.

\begin{acknowledgments}
Computations performed with PARI/GP (cypari2), SageMath, SnaPPy,
mpmath, and sympy. The Wright Brothers Research Program is
an independent theoretical physics initiative.
\end{acknowledgments}

\end{document}
